%% ---- Նախագծի տվյալները ----
\settopictitle{FPGA-ի հիման վրա սենսորային տվյալների հավաքագրման, մշակման և վերլուծության ինտեգրված ապարատածրագրային հարթակի մշակում}
\settopicsubtitle{FPGA-Based Sensor Data Acquisition, Processing and Analytics Platform}
\setversion{1.0}
\setdate{}
\setstatus{Օրինակ}

\setinstitution{[Համալսարանի անվանում]}
\setfaculty{[Ֆակուլտետ]}
\setcity{Երևան}
\setyear{2026}

\setstudentone{}
\setdeptone{Միկրոէլեկտրոնիկայի ամբիոն}
\setsupervisorone{[Ղեկավարի անուն, ազգանուն, աստիճան]}
\setstudenttwo{}
\setdepttwo{Ծրագրային / Համակարգչային ճարտարագիտության ամբիոն}
\setsupervisortwo{[Ղեկավարի անուն, ազգանուն, աստիճան]}
\setpartnerorg{}
\setmentor{}

\setworktype{Բակալավրի դիպլոմային նախագիծ}
\setdegreelevel{Բակալավր}
\setareas{Թվային նախագծում, FPGA, ներդրված համակարգեր, Web հարթակներ, տվյալների վերլուծություն}
\setkeywords{FPGA, Zynq UltraScale+, AXI, RTL, սենսոր-էմուլյացիա, runtime կոնֆիգուրացիա, տվյալների վերլուծություն}
\setslots{2 (մեկական յուրաքանչյուր ամբիոնից)}
\setduration{2 կիսամյակ (մոտ 30 շաբաթ)}
\setstartdate{Սեպտեմբեր, 2026}
\setdeadline{[ամսաթիվ]}
\setworklanguage{Հայերեն}
\settools{AMD AUP ZU3 (Zynq UltraScale+ MPSoC), Vivado, Vitis; սերվերային համակարգիչ, Python, SQL տվյալների բազա}
\setprerequisites{Թվային սխեմաներ և HDL (ուսանող 1); ծրագրավորում, տվյալների բազաներ, Web (ուսանող 2)}
